\documentclass[11pt,letterpaper]{article}

\usepackage[margin=0.82in,headheight=15pt,footskip=28pt]{geometry}
\usepackage{amsthm}
\usepackage{newtxtext,newtxmath}
\usepackage[T1]{fontenc}
\usepackage{microtype}
\usepackage{mathtools,bm}
\usepackage{booktabs,tabularx,array}
\usepackage{enumitem}
\usepackage{xcolor}
\usepackage[most]{tcolorbox}
\usepackage{titlesec}
\usepackage{fancyhdr}
\usepackage{hyperref}
\usepackage{lastpage}

\definecolor{onyx}{HTML}{17191C}
\definecolor{graphite}{HTML}{4B4F55}
\definecolor{rulegray}{HTML}{B8BBC0}
\definecolor{papergray}{HTML}{F3F3F1}
\definecolor{gold}{HTML}{8B6B2D}
\definecolor{deepblue}{HTML}{244B66}

\hypersetup{
  colorlinks=true,
  linkcolor=onyx,
  urlcolor=deepblue,
  pdftitle={A Complete 2x2 Operator Algebra},
  pdfauthor={Anthony Vito Coppa},
  pdfsubject={Recurrence, phase, and reversible projective transport in the full real 2x2 matrix algebra},
  pdfkeywords={matrix algebra, Fibonacci matrix, normal form, projective dynamics, dihedral group, Clifford algebra}
}

\pagestyle{fancy}
\fancyhf{}
\fancyhead[L]{\small A. V. Coppa}
\fancyhead[R]{\small A Complete $2\times2$ Operator Algebra}
\fancyfoot[L]{\small A. V. Coppa -- 5 August 2026}
\fancyfoot[R]{\small \thepage\ of \pageref{LastPage}}
\renewcommand{\headrulewidth}{0.35pt}
\renewcommand{\footrulewidth}{0pt}

\titleformat{\section}
  {\Large\bfseries\color{onyx}}
  {\thesection}{0.7em}{}
  [\vspace{0.2em}\color{rulegray}\titlerule]
\titleformat{\subsection}
  {\large\bfseries\color{onyx}}
  {\thesubsection}{0.65em}{}
\titlespacing*{\section}{0pt}{2.2ex plus .4ex minus .2ex}{1.2ex}
\titlespacing*{\subsection}{0pt}{1.8ex plus .3ex minus .2ex}{0.8ex}

\setlist[itemize]{leftmargin=1.4em,itemsep=0.22em,topsep=0.35em}
\setlist[enumerate]{leftmargin=1.55em,itemsep=0.22em,topsep=0.35em}
\setlength{\parindent}{0pt}
\setlength{\parskip}{0.52em}
\renewcommand{\arraystretch}{1.15}

\newtheoremstyle{kernelstyle}
  {0.7em}{0.7em}{\itshape}{}% above below body indent
  {\bfseries\color{onyx}}{.}{0.5em}{}
\theoremstyle{kernelstyle}
\newtheorem{theorem}{Theorem}[section]
\newtheorem{proposition}[theorem]{Proposition}
\newtheorem{corollary}[theorem]{Corollary}
\newtheorem{definition}[theorem]{Definition}
\newtheorem{lemma}[theorem]{Lemma}
\newtheorem{remark}[theorem]{Remark}

\newcommand{\R}{\mathbb{R}}
\newcommand{\Qfield}{\mathbb{Q}}
\newcommand{\Z}{\mathbb{Z}}
\newcommand{\N}{\mathbb{N}}
\newcommand{\PGL}{\operatorname{PGL}}
\newcommand{\GL}{\operatorname{GL}}
\newcommand{\End}{\operatorname{End}}
\newcommand{\Alg}{\operatorname{Alg}}
\newcommand{\Span}{\operatorname{span}}
\newcommand{\NF}{\operatorname{NF}}
\newcommand{\Eval}{\operatorname{Eval}}
\newcommand{\Fix}{\operatorname{Fix}}
\newcommand{\Eq}{\operatorname{Eq}}
\newcommand{\Dih}{\operatorname{Dih}}
\newcommand{\Cl}{\operatorname{Cl}}
\newcommand{\diag}{\operatorname{diag}}
\newcommand{\ad}{\operatorname{ad}}
\newcommand{\vecc}{\operatorname{vec}}
\newcommand{\id}{\operatorname{id}}

\newtcolorbox{abstractbox}{
  colback=papergray,
  colframe=rulegray,
  boxrule=0.5pt,
  arc=1.2pt,
  left=8pt,right=8pt,top=7pt,bottom=7pt
}
\newtcolorbox{boundarybox}{
  colback=white,
  colframe=gold,
  boxrule=0.65pt,
  arc=1.2pt,
  left=8pt,right=8pt,top=7pt,bottom=7pt,
  title={\bfseries Scope and evidence boundary},
  coltitle=onyx,
  fonttitle=\normalsize
}
\newtcolorbox{resultbox}{
  colback=papergray,
  colframe=onyx,
  boxrule=0.7pt,
  arc=1.2pt,
  left=9pt,right=9pt,top=8pt,bottom=8pt
}

\begin{document}
\hypersetup{pageanchor=false}

\begin{titlepage}
\thispagestyle{empty}
\vspace*{0.10\textheight}

{\fontsize{29}{34}\selectfont\bfseries\color{onyx} A Complete $2\times2$ Operator Algebra\par}
\vspace{0.7em}
{\fontsize{17}{22}\selectfont\bfseries Recurrence, Phase, and Reversible Projective Transport\par}

\vspace{2.1em}
\color{gold}\rule{0.62\textwidth}{1.1pt}\color{black}
\vspace{2.0em}

\begin{abstractbox}
\textbf{Abstract.}
Let
\[
Q=\begin{pmatrix}1&1\\1&0\end{pmatrix},
\qquad
C_1=\begin{pmatrix}0&-1\\1&0\end{pmatrix},
\qquad
\Omega=-C_1.
\]
The algebra generated by $Q$ and $C_1$ is the full real matrix algebra $M_2(\R)$.  The order defect
\[
G=Q^{\mathsf T}\Omega-\Omega Q=[C_1,Q]
\]
completes the basis $\{I,Q,\Omega,G\}$, so every real $2\times2$ operator has a unique four-coordinate normal form.  A finite multiplication table reduces arbitrary operator words exactly.  The same kernel carries a quarter-turn phase flow, Fibonacci recurrence, a symmetric signature-$(1,1)$ form whose null rays are the two eigendirections of $Q$, a normalized reflection, finite square symmetry, infinite-dihedral projective transport, an exact four-phase quotient, and bounded dominant-eigenvalue normalization.  Thus the kernel supplies infinite compositional syntax with finite exact semantics.  No occurrence ledger, covering-space return, finite eight-state carrier, or physical interpretation is used in the proofs.
\end{abstractbox}

\vfill
{\large Anthony Vito Coppa\par}
{\large Independent Research\par}
{\large Draft 0.2 -- 5 August 2026\par}
\vspace{0.6em}
{\small Jurisdiction: real linear algebra, projective geometry, elementary discrete groups, and the hyperbolic plane.\par}
\end{titlepage}
\hypersetup{pageanchor=true}
\pagenumbering{arabic}

\tableofcontents
\clearpage

\section{Purpose, jurisdiction, and main statement}

This paper isolates a complete planar operator algebra. Its burden is deliberately narrow: construct the exact local operator language, prove every reduction and transport statement in its proper carrier, and state the resulting data without importing bookkeeping, spin, finite relational addressing, gauge fields, or semantic interpretation.

\begin{boundarybox}
The following distinctions are part of the result.
\begin{itemize}
  \item A matrix, bilinear form, projective transformation, group element, spectral value, and normal-form coordinate occupy different mathematical offices even when they share a numerical representative.
  \item The names ``recurrence,'' ``phase,'' ``reflection,'' and ``transport'' refer only to the operators defined below.
  \item The constants $\Phi$, $\Psi$, and $\sqrt5$ enter through the characteristic polynomial, eigenstructure, and normalization of explicit matrices.  No external interpretation is required.
  \item Physical realization and retained occurrence lie outside the jurisdiction of this paper and are not premises here.
\end{itemize}
\end{boundarybox}

\begin{theorem}[K.0 -- Complete operator-algebra theorem]\label{thm:main}
There is a two-generator real operator kernel
\[
\mathcal K_2=(\R^2,Q,C_1)
\]
with the following properties.
\begin{enumerate}
  \item $\Alg_{\R}(Q,C_1)=M_2(\R)$, and two generators are minimal by cardinality.
  \item Every $A\in M_2(\R)$ has a unique normal form
  \[
  A=x_0I+x_1Q+x_2\Omega+x_3G.
  \]
  \item Every finite word in $Q$ and $C_1$ reduces exactly by a finite multiplication table.
  \item $Q$ carries the Fibonacci recurrence, while $C_1$ carries continuous quarter-turn phase.
  \item The order defect $G$ defines a signature-$(1,1)$ symmetric form whose null rays are exactly the two eigendirections of $Q$.
  \item Normalized reflections generate the square group $\Dih_4$, while projective transport and reflection generate $D_\infty$, with exact quotient
  \[
  D_\infty/\langle q^4\rangle\cong\Dih_4.
  \]
  \item Dominant-eigenvalue normalization keeps $\Phi^{-c}Q^c$ bounded and makes its limiting projector explicit.
\end{enumerate}
\end{theorem}

The proof is assembled in Sections~\ref{sec:primitive}--\ref{sec:bounded}.  The export interface appears in Section~\ref{sec:export}.

\section{Primitive carrier, recurrence, and phase}\label{sec:primitive}

\begin{definition}[K.1 -- Primitive operators]
On $V=\R^2$, define
\begin{equation}\label{eq:base}
Q=
\begin{pmatrix}
1&1\\
1&0
\end{pmatrix},
\qquad
C_1=
\begin{pmatrix}
0&-1\\
1&0
\end{pmatrix},
\qquad
\Omega=-C_1=
\begin{pmatrix}
0&1\\
-1&0
\end{pmatrix}.
\end{equation}
Let $I$ denote the identity matrix.  The continuous phase carriage is
\begin{equation}\label{eq:phase}
U(\theta)=e^{\theta C_1}=I\cos\theta+C_1\sin\theta\in SO(2).
\end{equation}
\end{definition}

Directly,
\begin{equation}\label{eq:quarter}
C_1^2=\Omega^2=-I,
\qquad
C_1^4=\Omega^4=I.
\end{equation}

\begin{proposition}[K.2 -- Quadratic recurrence]\label{prop:quadratic}
The matrix $Q$ satisfies
\begin{equation}\label{eq:qquad}
Q^2=Q+I,
\qquad
Q^{-1}=Q-I.
\end{equation}
\end{proposition}

\begin{proof}
The characteristic polynomial of $Q$ is $\lambda^2-\lambda-1$.  Cayley--Hamilton gives $Q^2-Q-I=0$.  Since $\det Q=-1$, $Q$ is invertible, and multiplication by $Q^{-1}$ yields $Q^{-1}=Q-I$.
\end{proof}

\begin{proposition}[K.3 -- Fibonacci powers]\label{prop:fibpowers}
Let $F_0=0$, $F_1=1$, and $F_{n+1}=F_n+F_{n-1}$.  For every $n\ge1$,
\begin{equation}\label{eq:fibpower}
Q^n=F_nQ+F_{n-1}I
=
\begin{pmatrix}
F_{n+1}&F_n\\
F_n&F_{n-1}
\end{pmatrix}.
\end{equation}
\end{proposition}

\begin{proof}
The assertion is immediate for $n=1$.  If $Q^n=F_nQ+F_{n-1}I$, then Proposition~\ref{prop:quadratic} gives
\[
Q^{n+1}=F_n(Q+I)+F_{n-1}Q=F_{n+1}Q+F_nI.
\]
\end{proof}

This is the standard Fibonacci $Q$-matrix identity; see, for example, \cite{HardyWright}.

\begin{proposition}[K.8 -- Finite phase difference]\label{prop:phase-difference}
For every $\theta\in\R$,
\begin{equation}\label{eq:phase-difference}
U(\theta)-I
=
2\sin\!\left(\frac\theta2\right)
U\!\left(\frac\theta2\right)C_1.
\end{equation}
Consequently,
\begin{equation}
\|(U(\theta)-I)v\|
=2\left|\sin\!\left(\frac\theta2\right)\right|\|v\|.
\end{equation}
\end{proposition}

\begin{proof}
Using $C_1^2=-I$ and the half-angle identities,
\[
2\sin\!\left(\frac\theta2\right)
\left(I\cos\!\frac\theta2+C_1\sin\!\frac\theta2\right)C_1
=(\cos\theta-1)I+C_1\sin\theta.
\]
The norm formula follows because $U(\theta/2)$ and $C_1$ are orthogonal.
\end{proof}

\section{Generated difference and complete matrix algebra}\label{sec:algebra}

\begin{definition}[K.4 -- Order defect]
Define
\begin{equation}\label{eq:gdef}
G:=Q^{\mathsf T}\Omega-\Omega Q.
\end{equation}
Since $Q^{\mathsf T}=Q$ and $C_1=-\Omega$,
\begin{equation}\label{eq:gmatrix}
G=[C_1,Q]=
\begin{pmatrix}
-2&1\\
1&2
\end{pmatrix}.
\end{equation}
Thus $G$ measures the difference between quarter-turn-then-recurrence and recurrence-then-quarter-turn.
\end{definition}

\begin{theorem}[K.5 -- Complete local algebra]\label{thm:basis}
The ordered set
\begin{equation}
\mathcal B=(I,Q,\Omega,G)
\end{equation}
is a basis of $M_2(\R)$.  Hence
\begin{equation}
\Alg_{\R}(Q,C_1)=M_2(\R).
\end{equation}
No single real $2\times2$ matrix can generate $M_2(\R)$ as a unital real algebra, so the two-generator presentation is minimal by cardinality.
\end{theorem}

\begin{proof}
With column-wise coordinate convention
\[
\vecc\!\begin{pmatrix}a&b\\c&d\end{pmatrix}=(a,b,c,d)^{\mathsf T},
\]
the four basis candidates form the matrix
\[
B=
\begin{pmatrix}
1&1&0&-2\\
0&1&1&1\\
0&1&-1&1\\
1&0&0&2
\end{pmatrix},
\qquad
\det B=-10\ne0.
\]
They therefore form a basis of the four-dimensional vector space $M_2(\R)$.  Since $I,Q,\Omega,G$ all belong to the algebra generated by $Q$ and $C_1$, that algebra is all of $M_2(\R)$.

For minimality, the unital algebra generated by one matrix $A$ is $\R[A]$.  Cayley--Hamilton reduces every polynomial in $A$ to $xI+yA$, so $\dim\R[A]\le2$; it is also commutative.  It cannot equal the four-dimensional noncommutative algebra $M_2(\R)$.
\end{proof}

\begin{corollary}[K.6 -- Unique normal form]\label{cor:normalform}
Every
\[
A=\begin{pmatrix}a&b\\c&d\end{pmatrix}\in M_2(\R)
\]
has the unique form
\begin{equation}\label{eq:nf}
A=x_0I+x_1Q+x_2\Omega+x_3G,
\end{equation}
where
\begin{align}
 x_0&=\frac{2a+3d-b-c}{5},\label{eq:x0}\\
 x_1&=\frac{a-d+2b+2c}{5},\label{eq:x1}\\
 x_2&=\frac{b-c}{2},\label{eq:x2}\\
 x_3&=\frac{-2a+2d+b+c}{10}.\label{eq:x3}
\end{align}
\end{corollary}

\begin{proof}
Expand the right side of Eq.~\eqref{eq:nf} and solve the resulting four linear equations.  Uniqueness follows from Theorem~\ref{thm:basis}.
\end{proof}

\begin{theorem}[K.7 -- Exact word reduction]\label{thm:word-reduction}
Every finite word in the alphabet $\{Q,C_1\}$ evaluates to an integer matrix and reduces uniquely to the normal form of Corollary~\ref{cor:normalform}.  Reduction is governed by the finite table
\begin{align}
 Q^2&=I+Q, & \Omega^2&=-I, & G^2&=5I,\label{eq:table1}\\
 Q\Omega&=\frac{\Omega+G}{2}, &
 \Omega Q&=\frac{\Omega-G}{2},\label{eq:table2}\\
 QG&=\frac{G+5\Omega}{2}, &
 GQ&=\frac{G-5\Omega}{2},\label{eq:table3}\\
 \Omega G&=2Q-I, &
 G\Omega&=I-2Q.\label{eq:table4}
\end{align}
If $A$ and $B$ have normal-form coordinate vectors $x=(x_0,x_1,x_2,x_3)$ and $y=(y_0,y_1,y_2,y_3)$, then $AB$ has coordinates $z=(z_0,z_1,z_2,z_3)$ given by
\begin{align}
 z_0={}&x_0y_0+x_1y_1-x_2y_2-x_2y_3+x_3y_2+5x_3y_3,\label{eq:z0}\\
 z_1={}&x_0y_1+x_1y_0+x_1y_1+2x_2y_3-2x_3y_2,\label{eq:z1}\\
 z_2={}&x_0y_2+x_2y_0+\frac{x_1y_2+x_2y_1}{2}
       +\frac{5(x_1y_3-x_3y_1)}{2},\label{eq:z2}\\
 z_3={}&x_0y_3+x_3y_0+\frac{x_1y_2+x_1y_3-x_2y_1+x_3y_1}{2}.
\label{eq:z3}
\end{align}
\end{theorem}

\begin{proof}
Equations~\eqref{eq:table1}--\eqref{eq:table4} follow by direct matrix multiplication.  Any product of basis elements reduces to the basis, so induction on word length gives exact reduction of every finite word.  Equations~\eqref{eq:z0}--\eqref{eq:z3} follow by bilinear expansion through the same table.  Uniqueness follows from Theorem~\ref{thm:basis}.
\end{proof}

\begin{resultbox}
Let $\mathcal W=\{Q,C_1\}^{*}$ be the free monoid of finite operator words.  The kernel supplies the exact pipeline
\[
\mathcal W
\xrightarrow{\ \Eval\ }
M_2(\Z)
\xrightarrow{\ \NF\ }
\Qfield^4.
\]
The unreduced word may be retained as a construction record while the four-coordinate normal form supplies finite exact semantics.
\end{resultbox}

\section{Spectral directions and the quadratic audit}\label{sec:spectral}

The characteristic roots of $Q$ are
\begin{equation}\label{eq:golden}
\Phi=\frac{1+\sqrt5}{2},
\qquad
\Psi=\frac{1-\sqrt5}{2}=-\Phi^{-1},
\end{equation}
with
\begin{equation}\label{eq:golden-relations}
\Phi+\Psi=1,
\qquad
\Phi\Psi=-1,
\qquad
\Phi-\Psi=\sqrt5.
\end{equation}
Define
\begin{equation}
 v_\Phi=\binom{\Phi}{1},
 \qquad
 v_\Psi=\binom{\Psi}{1}.
\end{equation}

\begin{theorem}[K.11 -- Spectral/null correspondence]\label{thm:golden}
The following statements hold.
\begin{enumerate}
  \item $Qv_\Phi=\Phi v_\Phi$ and $Qv_\Psi=\Psi v_\Psi$.
  \item The quadratic form
  \[
  \Delta_G(v)=v^{\mathsf T}Gv
  \]
  has signature $(1,1)$ and factors as
  \begin{equation}\label{eq:factor}
  \Delta_G(x,y)=2\left(y-\frac{x}{\Phi}\right)(y+\Phi x).
  \end{equation}
  \item Its projective null set is exactly
  \[
  \{[v_\Phi],[v_\Psi]\}.
  \]
  \item The oriented alternating-form separation is
  \begin{equation}\label{eq:separation}
  v_\Phi^{\mathsf T}\Omega v_\Psi=\sqrt5.
  \end{equation}
\end{enumerate}
\end{theorem}

\begin{proof}
The eigenvector identities follow by substitution.  Since $\det G=-5<0$, $G$ has signature $(1,1)$.  Direct expansion gives
\[
\Delta_G(x,y)=-2x^2+2xy+2y^2,
\]
and Eq.~\eqref{eq:factor} follows from $\Phi-\Phi^{-1}=1$.  The factors vanish precisely on $y=x/\Phi$ and $y=-\Phi x$, which are the same lines as $v_\Phi$ and $v_\Psi$.  Finally,
\[
\det(v_\Phi,v_\Psi)=\Phi-\Psi=\sqrt5.
\]
\end{proof}

\begin{proposition}[K.16 -- Maximal projective equalizer]\label{prop:equalizer}
Define the derived reflection matrix
\begin{equation}
M_0:=Q\Omega=
\begin{pmatrix}-1&1\\0&1\end{pmatrix}.
\end{equation}
Then $M_0^2=I$ and $Q=M_0C_1$.  On the affine chart $[s:1]$ of $\R P^1$, the projective involutions induced by $M_0$ and $C_1$ are
\begin{equation}
 m(s)=1-s,
 \qquad
 \sigma(s)=-\frac1s.
\end{equation}
Their equalizer is exactly $\{\Phi,\Psi\}$, and this is maximal for two distinct elements of $\PGL(2,\R)$.
\end{proposition}

\begin{proof}
The matrix statements are direct.  Equality $m(s)=\sigma(s)$ is equivalent to
\[
s^2-s-1=0,
\]
whose roots are $\Phi$ and $\Psi$.  For distinct projective maps $f,g$, their equalizer is the fixed set of $g^{-1}f$.  A nonidentity projective transformation of $\R P^1$ has at most two real eigenlines, hence at most two fixed points.
\end{proof}

\section{Normalized reflections and Clifford form}\label{sec:clifford}

We use the standard real Clifford-algebra convention for signature $(1,1)$; see \cite{LawsonMichelsohn}.

\begin{definition}
Define
\begin{equation}\label{eq:hands}
H:=H_{\sqrt5}=\frac{G}{\sqrt5},
\qquad
S:=S_Q=\frac{2Q-I}{\sqrt5}.
\end{equation}
\end{definition}

\begin{theorem}[K.12 -- Normalized involutions]\label{thm:involutions}
The matrices $H$ and $S$ are Euclidean reflections:
\begin{equation}\label{eq:reflections}
H^2=S^2=I,
\qquad
H^{\mathsf T}H=S^{\mathsf T}S=I,
\qquad
\det H=\det S=-1.
\end{equation}
Moreover,
\begin{equation}\label{eq:anti}
H\Omega=-\Omega H,
\qquad
SH=\Omega,
\qquad
HS=-\Omega,
\end{equation}
and $H$ exchanges the two spectral directions projectively:
\begin{equation}\label{eq:exchange}
Hv_\Phi=\Phi v_\Psi,
\qquad
Hv_\Psi=\Phi^{-1}v_\Phi.
\end{equation}
\end{theorem}

\begin{proof}
Direct multiplication gives
\[
G^2=5I,
\qquad
(2Q-I)^2=5I,
\qquad
2Q-I=\Omega G,
\qquad
G\Omega=-\Omega G.
\]
Equations~\eqref{eq:reflections}--\eqref{eq:anti} follow.  Equation~\eqref{eq:exchange} follows by substitution using the identities in Eq.~\eqref{eq:golden-relations}.
\end{proof}

\begin{proposition}[K.20 -- Clifford realization]\label{prop:clifford}
The assignment
\begin{equation}
 e_+\longmapsto H,
 \qquad
 e_-\longmapsto\Omega
\end{equation}
extends to an algebra isomorphism
\begin{equation}
\Cl_{1,1}(\R)\cong M_2(\R),
\end{equation}
where $e_+^2=+1$, $e_-^2=-1$, and $e_+e_-=-e_-e_+$.  The recurrence operator is recovered by
\begin{equation}\label{eq:qrecover}
Q=\frac{I+\Omega G}{2}
=\frac{I+\sqrt5\,\Omega H}{2}.
\end{equation}
\end{proposition}

\begin{proof}
The defining Clifford relations follow from Theorem~\ref{thm:involutions} and $\Omega^2=-I$.  Their image spans
\[
\{I,H,\Omega,H\Omega\}=M_2(\R)
\]
by Theorem~\ref{thm:basis}; both algebras have real dimension four, so the homomorphism is an isomorphism.  Equation~\eqref{eq:qrecover} follows from $\Omega G=2Q-I$.
\end{proof}

\section{Finite and infinite reversible transport}\label{sec:groups}

\begin{theorem}[K.13 -- Finite square group]\label{thm:d4}
The matrices $\Omega$ and $H$ satisfy
\begin{equation}
\Omega^4=I,
\qquad
H^2=I,
\qquad
H\Omega H=\Omega^{-1}.
\end{equation}
Hence
\begin{equation}
\langle\Omega,H\rangle\cong\Dih_4,
\end{equation}
where $\Dih_4$ denotes the order-eight symmetry group of the square.
\end{theorem}

\begin{proof}
The relations are immediate from Eqs.~\eqref{eq:quarter} and \eqref{eq:anti}.  The subgroup contains four distinct rotations $I,\Omega,-I,-\Omega$ and four reflections obtained by left multiplication by $H$, so it has order eight and the standard dihedral presentation.
\end{proof}

\begin{theorem}[K.14 -- Infinite-dihedral projective transport]\label{thm:dinf}
Let
\[
q=[Q],
\qquad
h=[H]
\]
in $\PGL(2,\R)$.  Then
\begin{equation}\label{eq:inverseconj}
h^2=1,
\qquad
hqh=q^{-1},
\end{equation}
and $q$ has infinite order.  Therefore
\begin{equation}
\langle q,h\rangle\cong D_\infty.
\end{equation}
\end{theorem}

\begin{proof}
Direct multiplication gives
\begin{equation}\label{eq:hqh}
HQH=-Q^{-1}.
\end{equation}
The scalar sign disappears in projective space, giving Eq.~\eqref{eq:inverseconj}.  The eigenvalue ratio of $Q$ is
\[
\frac{\Phi}{\Psi}=-\Phi^2,
\]
which is not a root of unity.  Hence no positive power of $Q$ is scalar, so $q$ has infinite order.  An involution conjugating an infinite-order generator to its inverse gives the standard infinite-dihedral presentation.
\end{proof}

\begin{theorem}[K.15 -- Exact four-phase quotient]\label{thm:quotient}
The assignment
\begin{equation}
\pi_4(q)=\Omega,
\qquad
\pi_4(h)=H
\end{equation}
defines a surjective homomorphism
\begin{equation}
\pi_4:D_\infty\twoheadrightarrow\Dih_4
\end{equation}
with kernel $\langle q^4\rangle$.  Therefore
\begin{equation}\label{eq:quotient}
D_\infty/\langle q^4\rangle\cong\Dih_4.
\end{equation}
\end{theorem}

\begin{proof}
The images satisfy the defining relations by Theorem~\ref{thm:d4}, so the map is well defined and surjective.  Every element of $D_\infty$ is $q^n$ or $hq^n$.  The first maps to the identity exactly when $4\mid n$.  The second maps to a reflection $H\Omega^n$, which cannot be the identity because its determinant is $-1$.  Hence the kernel is exactly $\langle q^4\rangle$, and the first isomorphism theorem gives Eq.~\eqref{eq:quotient}.
\end{proof}

\section{Projective and hyperbolic realization}\label{sec:hyperbolic}

We use the standard extended action of $\GL(2,\R)$ on the upper half-plane; see \cite{Beardon,Lang}.

On the affine real line, projective $Q$-transport is
\begin{equation}\label{eq:mobius}
T(s)=1+\frac1s.
\end{equation}
Its fixed points are $\Phi$ and $\Psi$.  Proposition~\ref{prop:fibpowers} gives
\begin{equation}\label{eq:orbit}
T^n(\infty)=\frac{F_{n+1}}{F_n}
\qquad(n\ge1).
\end{equation}

\begin{proposition}[K.17 -- Hyperbolic glide realization]\label{prop:glide}
Extend the action of $A=\begin{psmallmatrix}a&b\\c&d\end{psmallmatrix}\in\GL(2,\R)$ to the upper half-plane by
\begin{equation}
\rho_A(z)=
\begin{cases}
\dfrac{az+b}{cz+d},&\det A>0,\\[6pt]
\dfrac{a\overline z+b}{c\overline z+d},&\det A<0.
\end{cases}
\end{equation}
Then $\rho_Q$ is an orientation-reversing glide reflection along the geodesic with ideal endpoints $\Psi$ and $\Phi$.  Its glide displacement is $2\log\Phi$, and $\rho_{Q^2}$ is the corresponding hyperbolic translation of length $4\log\Phi$.
\end{proposition}

\begin{proof}
The boundary fixed points solve $s=T(s)$ and are therefore $\Phi$ and $\Psi$.  Since $\det Q=-1$, $\rho_Q$ reverses orientation; it has no fixed point in the upper half-plane and preserves the geodesic joining its boundary fixed points, so it is a glide reflection.  The matrix $Q^2\in SL(2,\R)$ has eigenvalues $\Phi^2$ and $\Phi^{-2}$ and trace $3$.  Its translation length is
\[
2\operatorname{arcosh}(3/2)=4\log\Phi.
\]
A glide reflection moves by half the translation length of its square.
\end{proof}

\section{Unattained limit and bounded normalized evolution}\label{sec:bounded}

\begin{proposition}[K.18 -- Unattained projective limit]\label{prop:horizon}
The finite rational orbit in Eq.~\eqref{eq:orbit} satisfies
\begin{equation}
\lim_{n\to\infty}\frac{F_{n+1}}{F_n}=\Phi,
\end{equation}
but
\begin{equation}
\frac{F_{n+1}}{F_n}\ne\Phi
\end{equation}
for every finite $n\ge1$.
\end{proposition}

\begin{proof}
Convergence follows from Binet's formula.  Every finite ratio is rational, while $\Phi$ is irrational.
\end{proof}

Define the spectral projectors
\begin{equation}\label{eq:projectors}
P_\Phi=\frac{Q-\Psi I}{\sqrt5},
\qquad
P_\Psi=\frac{\Phi I-Q}{\sqrt5}.
\end{equation}
They satisfy
\begin{equation}
P_\Phi+P_\Psi=I,
\qquad
P_\Phi P_\Psi=0,
\qquad
QP_\Phi=\Phi P_\Phi,
\qquad
QP_\Psi=\Psi P_\Psi.
\end{equation}

\begin{theorem}[K.19 -- Bounded normalized recurrence]\label{thm:normalized}
For every $c\in\N$,
\begin{equation}\label{eq:normalized}
\Phi^{-c}Q^c
=P_\Phi+(-\Phi^{-2})^cP_\Psi.
\end{equation}
Hence the normalized sequence is bounded and converges exponentially to $P_\Phi$.
\end{theorem}

\begin{proof}
The spectral decomposition is
\[
Q^c=\Phi^cP_\Phi+\Psi^cP_\Psi.
\]
Divide by $\Phi^c$ and use $\Psi/\Phi=-\Phi^{-2}$.  Since $0<\Phi^{-2}<1$, the second term tends exponentially to zero.
\end{proof}

\section{Proved data and strict scope}\label{sec:export}

The proved data of this paper may be summarized by the tuple
\begin{equation}\label{eq:export}
\mathsf{KernelData}
=
\bigl(
V,Q,C_1,\Omega,G,H,S,\mathcal B,\NF,U,q,h,\pi_4,P_\Phi,P_\Psi
\bigr),
\end{equation}
together with the proofs attached to Theorems~\ref{thm:basis}, \ref{thm:word-reduction}, \ref{thm:golden}, \ref{thm:d4}, \ref{thm:dinf}, \ref{thm:quotient}, and \ref{thm:normalized}.

Equation~\eqref{eq:export} is a compact register of the objects established above, not an additional identification among them. In particular:
\begin{itemize}
  \item $Q^2$ restores the two bilinear forms under pullback; it is not a stored occurrence.
  \item $\Delta_G(v)=0$ identifies a null direction of a declared symmetric form; it is not a vanished vector or absent event.
  \item $D_\infty\twoheadrightarrow\Dih_4$ is a quotient of addressed projective transport; it does not itself supply chronology.
  \item $\Phi^{-c}Q^c$ is bounded after normalization; the integer lift count $c$ remains external to the normalized matrix value.
  \item The $137$-address register, eight-channel carrier, and twenty-eight pair relations are not forced by this kernel and therefore do not appear in its theorem.
  \item No claim about spin, gauge curvature, material dynamics, atomic structure, or semantic interpretation follows without a separately defined adapter or realization theorem.
\end{itemize}

\begin{proof}[Proof of Theorem~\ref{thm:main}]
Items 1--3 are Theorems~\ref{thm:basis} and \ref{thm:word-reduction}.  Item 4 is Propositions~\ref{prop:fibpowers} and \ref{prop:phase-difference}.  Item 5 is Theorem~\ref{thm:golden}.  Item 6 is Theorems~\ref{thm:d4}, \ref{thm:dinf}, and \ref{thm:quotient}.  Item 7 is Theorem~\ref{thm:normalized}.  Every object and map occurs in the export tuple with a recoverable carrier and codomain.
\end{proof}

\section{Computational verification}

The companion regression script checks all displayed finite matrix identities, the basis determinant, normal-form extraction, multiplication table, symbolic coordinate product, spectral projectors, and exhaustive word reduction through a declared finite depth. These computations are regression gates for the formulas; the universal proofs remain the arguments above.

\begin{thebibliography}{9}
\small

\bibitem{Beardon}
A.~F. Beardon,
\emph{The Geometry of Discrete Groups},
Springer, New York, 1983.

\bibitem{HardyWright}
G.~H. Hardy and E.~M. Wright,
\emph{An Introduction to the Theory of Numbers},
6th ed., Oxford University Press, 2008.

\bibitem{Lang}
S. Lang,
\emph{$SL_2(\mathbb R)$},
Springer, New York, 1985.

\bibitem{LawsonMichelsohn}
H.~B. Lawson, Jr. and M.-L. Michelsohn,
\emph{Spin Geometry},
Princeton University Press, 1989.

\end{thebibliography}

\end{document}
